% Source: https://artofproblemsolving.com/community/c3h1741603_find_the_third_side_of_the_triangle
% Stable identity: aops-topic-1741603-q01
Let $AB=5$ and $AC=6$, with circumradius $R=5$. By the extended law of sines, $\sin C=\dfrac{AB}{2R}=\dfrac{1}{2}$ and $\sin B=\dfrac{AC}{2R}=\dfrac{3}{5}$. So $C=30^\circ$ or $150^\circ$, and $\cos B=\pm\dfrac{4}{5}$.

Since $C$ is opposite the shorter side $AB<AC$, $C<B$, so $C=30^\circ$ ($C=150^\circ$ would leave $B<30^\circ$). Two cases remain, $B$ acute or obtuse, and both occur:
\begin{center}

\begin{asy}
usepackage("amsmath");
size(4.2cm);
defaultpen(fontsize(10pt));
pair O=(0,0);pair A=5*dir(90),B=5*dir(150),C=5*dir(90-2*aSin(0.6));
draw(circle(O,5),gray);draw(A--B--C--cycle,linewidth(0.9));
dot(O);label("$O$",O,S);
label("$A$",A,N);label("$B$",B,W);label("$C$",C,E);
label("$5$",(A+B)/2+1.15*dir(120));label("$6$",(A+C)/2,dir(50));
label("\mbox{}",(min(currentpicture,user=true)+max(currentpicture,user=true))/2,black); // keeps the text after the figure black when LaTeX typesets the labels inline
\end{asy}
\qquad

\begin{asy}
usepackage("amsmath");
size(4.2cm);
defaultpen(fontsize(10pt));
pair O=(0,0);pair A=5*dir(90),B=5*dir(150),C=5*dir(90+2*aSin(0.6));
draw(circle(O,5),gray);draw(A--B--C--cycle,linewidth(0.9));
dot(O);label("$O$",O,S);
label("$A$",A,N);label("$B$",B,NW);label("$C$",C,W);
label("$5$",(A+B)/2+1.15*dir(120));label("$6$",(A+C)/2,dir(-30));
label("\mbox{}",(min(currentpicture,user=true)+max(currentpicture,user=true))/2,black); // keeps the text after the figure black when LaTeX typesets the labels inline
\end{asy}

\end{center}
The third side is $BC=2R\sin A=10\sin(B+C)=10(\sin B\cos C+\cos B\sin C)=10\left(\dfrac{3}{5}\cdot\dfrac{\sqrt3}{2}\pm\dfrac{4}{5}\cdot\dfrac{1}{2}\right)$:
\begin{itemize}
\item $B$ acute (left): $BC=3\sqrt3+4\approx9.20$;
\item $B$ obtuse (right): $BC=3\sqrt3-4\approx1.20$.
\end{itemize}
So the third side is $\boxed{3\sqrt3+4\text{ or }3\sqrt3-4}$.
